\section{前向过程与并行执行}
\label{sec:algorithm}

初版路线规划完全确定，并在任何Mamba扫描之前结束。完整前向过程如下：

\begin{figure}[t]
\refstepcounter{algorithm}\label{alg:topocontour}
\small
\noindent\textbf{算法 \thealgorithm：\method{}前向过程}
\vspace{2pt}\hrule\vspace{3pt}
\begin{enumerate}[label=\arabic*.,leftmargin=1.5em,itemsep=1.5pt]
    \item 归一化原图RGB，计算确定性显著性场，完成平滑并生成一张Sobel场。
    \item 在补边后的分片线性显著性场上建立完整等高线树；提取保留轮廓、过滤小圈、插入树弧，并收缩没有采样轮廓的路径。
    \item 每个圈统一定向并按周长确定点数。均匀试探点读取带符号的法线Sobel响应：平均符号决定整圈低值普通侧，平滑强度决定密度，最大强度位置决定 $p_0$。
    \item 从 $p_0$ 开始按加权弧长重分配固定点数。只在最终点上查询普通侧的首次碰撞。
    \item 沿所有普通法线读取原图RGB与路线属性，从碰撞点扫描到带Tag的轮廓点；构造轮廓点Token，再从 $p_0$ 开始按两种方向扫描到带Tag返回点，得到 $y_C$。
    \item 每个终止圈还扫描相反侧法线，与第一次闭环的逐点输出融合，再重复两条从 $p_0$ 到带Tag返回点的闭环，得到 $y_C^{\rm term}$。
    \item 从低到高处理稀疏树：Tree Mamba扫描非分叉树段；分裂复制带Tag的 $y$，汇合进行支撑量感知的无序融合。
    \item 根据终止圈实际读取的两侧法线计算叶子显著性，并按弧长支撑量加权；叶子从低到高排序后交给最终Merge Mamba。
    \item 若没有有效圈，则投影原图RGB各通道的均值和方差；输出分类结果。
\end{enumerate}
\vspace{2pt}\hrule
\end{figure}

不同长度的法线、闭环和树段采用动态长度桶分批。填充位置位于真实带Tag终点之后并被屏蔽。所有法线与局部轮廓编码器相互独立，可以跨层并行；树上同一就绪前沿的树段也可并行。必须保持顺序的部分只有单条法线或闭环内部、低到高的树边以及最终叶子序列。长圈使用更大的桶或不丢点的精确分块，不允许截断或几何拉伸。

跨树段只传递带Tag的 $y$。分裂复制一个向量，而不是复制参数、隐藏状态或循环缓存。初版中，原图属性的逐点投影和所有Mamba模块可以训练；传统显著性算法、Sobel决策、轮廓提取和硬树结构都是停止梯度的路线构建。未来若加入显著性CNN，需要单独的软训练路径。
